\chapter{球座標系}
    \section{計量係数}
        \begin{shadebox}
            \[h_r = 1, \hth = r, \hph = r\sinth\]
        \end{shadebox}
        \begin{proof}
            \quad\par
            位置ベクトルは
            \[\vr = \dthv{r\sinth\cosph}{r\sinth\sinph}{r\costh}\]
            よって
            \begin{align*}
                h_r &\coloneqq \norm{\partDeriv{\vr}{r}{}} = \norm{\dthv{\sinth\cosph}{\sinth\sinph}{\costh}} = 1\\
                \hth &\coloneqq \norm{\partDeriv{\vr}{\theta}{}} = \norm{\dthv{r\costh\cosph}{r\costh\sinph}{-r\sinth}} = r\\
                \hph &\coloneqq \norm{\partDeriv{\vr}{\varphi}{}} = \norm{\dthv{-r\sinth\sinph}{r\sinth\cosph}{0}} = r\sinth
            \end{align*}
        \end{proof}
    \section{単位ベクトルに関する3次元Euclid空間との関係}
        \begin{shadebox}
            \begin{align*}
                \ir &= \vix\sinth\cosph + \viy\sinth\sinph + \viz\costh\\
                \ith &= \vix\costh\cosph + \viy\costh\sinph - \viz\sinth\\
                \iph &= -\vix\sinph + \viy\cosph
            \end{align*}
            \begin{align*}
                \vix &= \ir\sinth\cosph + \ith\costh\cosph -\iph\sinph\\
                \viy &= \ir\sinth\sinph + \ith\costh\sinph + \iph\cosph\\
                \viz &= \ir\costh - \ith\sinth
            \end{align*}
        \end{shadebox}
        \noindent(手間が大きい証明)
        \begin{proof}
            \quad\par
            位置ベクトルは
            \[\vr = \dthv{r\sinth\cosph}{r\sinth\sinph}{r\costh}\]
            よって
            \begin{align}
                \ir &\coloneqq \frac{1}{h_r}\partDeriv{\vr}{r}{} = \dthv{\sinth\cosph}{\sinth\sinph}{\costh} = \vix\sinth\cosph + \viy\sinth\sinph + \viz\costh \\
                \ith &\coloneqq \frac{1}{\hth}\partDeriv{\vr}{\theta}{} = \frac{1}{r}\dthv{r\costh\cosph}{r\costh\sinph}{-r\sinth} = \vix\costh\cosph + \viy\costh\sinph - \viz\sinth \\
                \iph &\coloneqq \frac{1}{\hph}\partDeriv{\vr}{\varphi}{} = \frac{1}{r\sinth}\dthv{-r\sinth\sinph}{r\sinth\cosph}{0} = -\vix\sinph + \viy\cosph
            \end{align}
            これを逆に解いてみる。\\
            $(1)\times\costh-(2)\times\sinth$より$\viz = \ir\costh - \ith\sinth$\\
            $(1)\times\sinth+(2)\times\costh$より
            \begin{equation}
                \ir\sinth+\ith\costh = \vix\cosph+\viy\sinph
            \end{equation}
            $(4)\times\cosph-(3)\times\sinph$より$\vix = \ir\sinth\cosph + \ith\costh\cosph -\iph\sinph$\\
            $(3)\times\cosph+(4)\times\sinph$より$\viy = \ir\sinth\sinph + \ith\costh\sinph + \iph\cosph$
        \end{proof}
        \noindent(手間が少ない証明)
        \begin{proof}
            \quad\par
            $\ir,\ith,\iph$の導出までは上の証明と同じ。
            $\ir,\ith,\iph$が直交系を成すことを利用し、$\vix = (\vix\cdot\ir)\ir + (\vix\cdot\ith)\ith + (\vix\cdot\iph)\iph,\quad \viy = (\viy\cdot\ir)\ir + (\viy\cdot\ith)\ith + (\viy\cdot\iph)\iph,\quad \viz = (\viz\cdot\ir)\ir + (\viz\cdot\ith)\ith + (\viz\cdot\iph)\iph$として導出できる。
        \end{proof}
    % \section{Laplacian}
    % 	\begin{shadebox}
    % 		球座標系に於けるLaplace演算子は次式である。
    % 		\[ \Delta = \frac{1}{r^2}\frac{\partial}{\partial r}r^2\frac{\partial}{\partial r} + \frac{1}{r^2}\frac{\partial^2}{\partial\theta^2} + \frac{\cos\theta}{r^2\sin\theta}\frac{\partial}{\partial\theta} + \frac{1}{r^2\sin^2\theta}\frac{\partial^2}{\partial\varphi^2} \]
    % 	\end{shadebox}
    % 	\begin{proof}
    % 		\quad\par
    % 		直角座標を引数に取る関数$\phi:\realNumbers^3\to\complexNumbers$は$\mathrm{C}^2$級であるとする。
    % 		球座標から直角座標への変換関数を$f:[r,\theta,\varphi]^\top\in\realNumbers^3\mapsto[r\sin\theta\cos\varphi,r\sin\theta\sin\varphi,r\cos\theta]^\top\in\realNumbers^3$とする。
    % 		$\hat{\phi}(\bm{u}) := \phi(f(\bm{u}))$とする。
    % 		$\bm{r} := [x,y,z]^\top\in\realNumbers^3$とする。
    % 		\begin{align*}
    % 			\Delta_\bm{r}\phi(\bm{r}) &= \frac{\partial^2 \phi}{\partial x^2} + \frac{\partial^2 \phi}{\partial y^2} + \frac{\partial^2 \phi}{\partial z^2} \\
    % 			\begin{split}
    % 				&= \frac{\partial }{\partial x}\left(\frac{\partial \hat{\phi}}{\partial r}\frac{\partial r}{\partial x} + \frac{\partial \hat{\phi}}{\partial \theta}\frac{\partial \theta}{\partial x} + \frac{\partial \hat{\phi}}{\partial \varphi}\frac{\partial \varphi}{\partial x}\right) + \frac{\partial }{\partial y}\left(\frac{\partial \hat{\phi}}{\partial r}\frac{\partial r}{\partial y} + \frac{\partial \hat{\phi}}{\partial \theta}\frac{\partial \theta}{\partial y} + \frac{\partial \hat{\phi}}{\partial \varphi}\frac{\partial \varphi}{\partial y}\right) \\
    % 				&\phantom{=} + \frac{\partial }{\partial z}\left(\frac{\partial \hat{\phi}}{\partial r}\frac{\partial r}{\partial z} + \frac{\partial \hat{\phi}}{\partial \theta}\frac{\partial \theta}{\partial z} + \frac{\partial \hat{\phi}}{\partial \varphi}\frac{\partial \varphi}{\partial z}\right)
    % 			\end{split} \tag{1}
    % 		\end{align*}
    % 		右辺第1項を展開すると次式を得る。
    % 		\begin{align*}
    % 			&\phantom{=} \frac{\partial^2\phi}{\partial r^2}\left(\frac{\partial r}{\partial x}\right)^2 + \frac{\partial \hat{\phi}}{\partial r}\frac{\partial^2 r}{\partial x^2} + \frac{\partial^2\phi}{\partial \theta^2}\left(\frac{\partial \theta}{\partial x}\right)^2 + \frac{\partial \hat{\phi}}{\partial \theta}\frac{\partial^2 \theta}{\partial x^2} + \frac{\partial^2\phi}{\partial \varphi^2}\left(\frac{\partial \varphi}{\partial x}\right)^2 + \frac{\partial \hat{\phi}}{\partial \varphi}\frac{\partial^2 \varphi}{\partial x^2} \\
    % 			&\phantom{=} + 2\frac{\partial^2 \hat{\phi}}{\partial \theta \partial r}\frac{\partial \theta}{\partial x}\frac{\partial r}{\partial x} + 2\frac{\partial^2 \hat{\phi}}{\partial \varphi \partial \theta}\frac{\partial \varphi}{\partial x}\frac{\partial \theta}{\partial x} + 2\frac{\partial^2 \hat{\phi}}{\partial r \partial \varphi}\frac{\partial r}{\partial x}\frac{\partial \varphi}{\partial x}
    % 		\end{align*}
    % 		式(1)の右辺第2,3項はそれぞれ上式の$x$を$y,z$で置き換えたものになる。
    % 		以上より式(1)は次式になる。
    % 		\begin{align*}
    % 			(1) &= \frac{\partial^2 \hat{\phi}}{\partial r^2}\left[\left(\frac{\partial r}{\partial x}\right)^2 + \left(\frac{\partial r}{\partial y}\right)^2 + \left(\frac{\partial r}{\partial z}\right)^2\right] + \frac{\partial \hat{\phi}}{\partial r}\left(\frac{\partial^2 r}{\partial x^2} + \frac{\partial^2 r}{\partial y^2} + \frac{\partial^2 r}{\partial z^2}\right) \tag{2} \\
    % 			&\phantom{=} + \frac{\partial^2 \hat{\phi}}{\partial \theta^2}\left[\left(\frac{\partial \theta}{\partial x}\right)^2 + \left(\frac{\partial \theta}{\partial y}\right)^2 + \left(\frac{\partial \theta}{\partial z}\right)^2\right] + \frac{\partial \hat{\phi}}{\partial \theta}\left(\frac{\partial^2 \theta}{\partial x^2} + \frac{\partial^2 \theta}{\partial y^2} + \frac{\partial^2 \theta}{\partial z^2}\right) \tag{3} \\
    % 			&\phantom{=} + \frac{\partial^2 \hat{\phi}}{\partial \varphi^2}\left[\left(\frac{\partial \varphi}{\partial x}\right)^2 + \left(\frac{\partial \varphi}{\partial y}\right)^2 + \left(\frac{\partial \varphi}{\partial z}\right)^2\right] + \frac{\partial \hat{\phi}}{\partial \varphi}\left(\frac{\partial^2 \varphi}{\partial x^2} + \frac{\partial^2 \varphi}{\partial y^2} + \frac{\partial^2 \varphi}{\partial z^2}\right) \tag{4} \\
    % 			&\phantom{=} + 2\frac{\partial^2 \hat{\phi}}{\partial \theta \partial r}\left(\frac{\partial \theta}{\partial x}\frac{\partial r}{\partial x} + \frac{\partial \theta}{\partial y}\frac{\partial r}{\partial y} + \frac{\partial \theta}{\partial z}\frac{\partial r}{\partial z}\right) \tag{5} \\
    % 			&\phantom{=} + 2\frac{\partial^2 \hat{\phi}}{\partial \varphi \partial \theta}\left(\frac{\partial \varphi}{\partial x}\frac{\partial \theta}{\partial x} + \frac{\partial \varphi}{\partial y}\frac{\partial \theta}{\partial y} + \frac{\partial \varphi}{\partial z}\frac{\partial \theta}{\partial z}\right) \tag{6} \\
    % 			&\phantom{=} + 2\frac{\partial^2 \hat{\phi}}{\partial r \partial \varphi}\left(\frac{\partial r}{\partial x}\frac{\partial \varphi}{\partial x} + \frac{\partial r}{\partial y}\frac{\partial \varphi}{\partial y} + \frac{\partial r}{\partial z}\frac{\partial \varphi}{\partial z}\right) \tag{7}
    % 		\end{align*}
    % 		まず式(5),(6),(7)が0であることを示す。
    % 		$J$を$f$のJacobi行列とすると、式(5)の括弧の内側は$J^{-1}$の1行目と2行目の内積である。
    % 		球座標系は直交座標系であるため$J$の列は直交系を成す。
    % 		これと\ref{行列の列(行)の直交性と逆行列の行(列)の直交性}より$J^{-1}$の行は直交系を成すから、式(5)の括弧の内側は0である。
    % 		同様に式(6),(7)も0である。
    % 		\par
    % 		式(2),(3),(4)を評価する。
    % 		$r = \sqrt{x^2+y^2+z^2},\;\theta = \acos{z/r} = \acos{z/\sqrt{x^2+y^2+z^2}},\;\phi = \atanEx{x}{y}$
    % 		であるから次式が成り立つ。
    % 		\begin{gather*}
    % 			\frac{\partial r}{\partial x} = \sin\theta\cos\varphi, \quad \frac{\partial r}{\partial y} = \sin\theta\sin\varphi, \quad \frac{\partial r}{\partial z} = \cos\theta \\
    % 			\frac{\partial \theta}{\partial x} = \frac{1}{r}\cos\theta\cos\varphi, \quad \frac{\partial \theta}{\partial y} = \frac{1}{r}\cos\theta\sin\varphi, \quad \frac{\partial \theta}{\partial z} = -\frac{1}{r}\sin\theta \\
    % 			\frac{\partial \varphi}{\partial x} = -\frac{\sin\varphi}{r\sin\theta}, \quad \frac{\partial \varphi}{\partial y} = \frac{\cos\varphi}{r\sin\theta}, \quad \frac{\partial \varphi}{\partial z} = 0 \\
    % 			\frac{\partial^2 r}{\partial x^2} = \frac{1}{r}\left(\cos^2\theta\cos^2\varphi + \sin^2\varphi\right), \quad \frac{\partial^2 r}{\partial y^2} = \frac{1}{r}\left(\cos^2\theta\sin^2\varphi + \cos^2\varphi\right), \quad \frac{\partial^2 r}{\partial z^2} = \frac{1}{r}\sin^2\theta \\
    % 			\frac{\partial^2 \theta}{\partial x^2} = \frac{1}{r^2\sin\theta}\left(-2\sin^2\theta\cos\theta\cos^2\varphi + \cos\theta\sin^2\varphi\right) \\
    % 			\frac{\partial^2 \theta}{\partial y^2} = \frac{1}{r^2\sin\theta}\left(-2\sin^2\theta\cos\theta\sin^2\varphi + \cos\theta\cos^2\varphi\right) \\
    % 			\frac{\partial^2 \theta}{\partial z^2} = \frac{2}{r^2}\sin\theta\cos\theta \\
    % 			\frac{\partial^2 \varphi}{\partial x^2} = \frac{2\sin\varphi\cos\varphi}{r^2\sin^2\theta}, \quad \frac{\partial^2 \varphi}{\partial y^2} = -\frac{2\sin\varphi\cos\varphi}{r^2\sin^2\theta}
    % 		\end{gather*}
    % 		これらを用いて
    % 		\begin{align*}
    % 			(2) &= \frac{\partial^2 \hat{\phi}}{\partial r^2} + \frac{2}{r}\frac{\partial \hat{\phi}}{\partial r} = \frac{1}{r^2}\frac{\partial}{\partial r}r^2\frac{\partial \hat{\phi}}{\partial r} \\
    % 			(3) &= \frac{1}{r^2}\frac{\partial^2 \hat{\phi}}{\partial \theta^2} + \frac{\cos\theta}{r^2\sin\theta}\frac{\partial \hat{\phi}}{\partial \theta} \\
    % 			(4) &= \frac{1}{r^2\sin^2\theta}\frac{\partial^2 \hat{\phi}}{\partial \varphi^2}
    % 		\end{align*}
    % 	\end{proof}